%%%%%%%%%%%%%%%%%%%%%%%%%%%%%%%%%%%%%%%%%
% Medium Length Professional CV
% LaTeX Template
% Version 2.0 (8/5/13)
%
% This template has been downloaded from:
% http://www.LaTeXTemplates.com
%
% Original author:
% Trey Hunner (http://www.treyhunner.com/)
%
% Important note:
% This template requires the resume.cls file to be in the same directory as the
% .tex file. The resume.cls file provides the resume style used for structuring the
% document.
%
%%%%%%%%%%%%%%%%%%%%%%%%%%%%%%%%%%%%%%%%%

%----------------------------------------------------------------------------------------
%   PACKAGES AND OTHER DOCUMENT CONFIGURATIONS
%----------------------------------------------------------------------------------------

\documentclass{resume} % Use the custom resume.cls style
\usepackage[left=0.7in,top=0.7in,right=0.7in,bottom=0.7in]{geometry} % Document margins



%----------------------------------------------------------------------------------------
%   DOCUMENT START
%----------------------------------------------------------------------------------------
\name{Zhaoxuan (Hins) Hu} % Your name
\address{(217)-693-2283 \\ zh223@cornell.edu \\ GitHub: \href{https://github.com/Hins-Hu/}{hins-hu} \\ LinkedIn: \href{https://www.linkedin.com/in/hins-hu/}{hins-hu}}

% \href{www.github.com/hins-hu}{GitHub} \\ \href{www.linkedin.com/in/hins-hu}{LinkedIn}} %Your phone number and email

\begin{document}
\title{Hins Resume}
%----------------------------------------------------------------------------------------
%   EDUCATION SECTION
%----------------------------------------------------------------------------------------
% \begin{rSection}{Availability}
% I am available \textbf{from May 7, 2024 to August 31, 2024}, which covers all summer internship cycles
% \end{rSection}

% \begin{rSection}{Availability}
% \vspace{0.1cm}
% I am available \textbf{from May 7, 2024 to August 31, 2024}, which covers all summer internship cycles
% \end{rSection}

\vspace{0.2cm}
\begin{rSection}{Education}
\vspace{0.2cm}
\begin{rSubsection}{Cornell University}{Aug. 2021 - May 2026 (Expected)}{Ph.D. in Systems Engineering}{Ithaca, NY}
\item Minor in Computer Science (AI) and Operations Research (Mathematical Programming)
\item Focus on applying recent advancements (e.g., deep RL and LLM) in machine learning to solve (sequential) hard combinatorial optimization problems.
\end{rSubsection}
% \vspace{-0.2cm}
\begin{rSubsection}{University of Illinois at Urbana-Champaign}{Aug. 2019 - May 2021}{M.Sc. in Systems Engineering}{Urbana, IL}
\item[]
\vspace{-0.7cm}
\end{rSubsection}
% \vspace{-0.2cm}
\begin{rSubsection}{South China University of Technology}{Aug. 2015 - June 2019}{B.Eng. in Engineering (First Class Honor)}{Guangzhou, China}
% \item Exchange to Western University (Canada) in 2018-2019
\item[]
\vspace{-0.7cm}
\end{rSubsection}
\end{rSection}


\begin{rSection}{Work Experience}
\vspace{0.2cm}

\begin{rSubsection}{PhD Software Engineer Intern}{August 2025 - December 2025}{Google}{Mountain View, CA}
\item Host Team: Google Cloud - AI Flow 
\item Built core AI workflows to enable seamless LLM adoption in downstream products and surfaces.
\item Focused on the supervised fine-tuning and knowledge distillation of hybrid encoder-decoder Gemini models on text related-tasks. 
\end{rSubsection}

\vspace{0.1cm}

\begin{rSubsection}{PhD Machine Learning Engineer Intern}{May 2025 - August 2025}{Meta}{Bellevue, WA}
\item Host Team: Modern Recommendation Systems (MRS) - AI Knowledge
\item Contributed to the cross-surface HSTU model for generating embeddings that capture long-context user-item interactions, leading to a 0.2\% (Normalized Entropy) NE gain in top-line engagement metrics of the downstream ranking system, such as comment, like, and share.
\item Extended the module to integrate State Space Models (i.e., Mamba) for enhanced sequence modeling.
\end{rSubsection}
% \begin{rSubsection}{Pilot test of AI-powered micro-transit systems}{Jan. 2023 - May 2024}{Software Engineer, Remote}{Chattanooga Transit Agency}
% \item Contributed to the vehicle routing algorithms leveraging open-source engines like OSRM and Valhalla.
% \item Built tools to analyze mobility-energy efficiency of the deployed micro-transit system in target areas.
% \item Experienced in handling millions of ride requests and high-dimensional geospatial data.
% \end{rSubsection}
\end{rSection}

% \vspace{0.2cm}
\begin{rSection}{Selected Projects}
\vspace{0.3cm}
% \begin{rSubsection}{Learning for branch-and-cut}{Apr. 2023 - Present}{}{}
% \item Explore learning strategies to assist the tuning process of hyperparameters such as node selection, variable selection, and cut selection that are used in the classical branch-and-cut algorithm to solve a variety of combinatorial optimization problems
% \end{rSubsection}
\begin{rSubsection}{LLM Benchmark for Heuristics in Combinatorial Optimization}{Nov. 2024 - Oct. 2025}{}{}
\item Contributed to an agentic benchmark for LLM-crafted heuristics in real-world NP-hard combinatorial optimization problems across various domains and problem types.
\item Introducing Quality-Yield Index (QYI) – a unified metric for solution feasibility and quality.
\item Revealing that even SOTA models (e.g., Gemini-Pro) achieve less than 60\% of expert performance.
\item Refer to the project \href{https://cornell-zhang.github.io/heurigym/}{website} and \href{https://github.com/cornell-zhang/heurigym}{GitHub Repo} for more details.
\end{rSubsection}

\begin{rSubsection}{Learning algorithms for reliable online shortest path problems}{Oct. 2025 – Present}{}{}
\item Investigate the novel reliable online shortest path problem that accounts for both mean and standard deviation in path travel times.
\item Developed a provably efficient Lower Confidence Bound (LCB) algorithm for learning routing strategies within a combinatorial multi-armed bandit framework.    
\end{rSubsection}

\begin{rSubsection}{End-to-end deep learning for combinatorially constrained MDP}{May 2025 - Present}{}{}
\item Developed an end-to-end constrained optimization learning framework to solve the Markov decision process (MDP) with combinatorially constrained action space and tested it in ride-sharing systems.
\item Designed a GPT-based sequence model to learn massive offline trajectories of on-demand ride-sharing fleets given by classical myopic optimization algorithms.
\end{rSubsection}


\begin{rSubsection}{GNN for the algorithmic speed-up in ride-pool matching}{Nov. 2023 - Nov. 2024}{}{}
\item Developed a GNN-based supervised learning model to prune the request-trip-vehicle graph used for generating candidate shared trips.
\item The learning-augmented framework outperforms the hand-crafted heuristic counterpart by a 5x speed-up while maintaining a negligible impact on the solution quality.
\end{rSubsection}

% \begin{rSubsection}{Reinforcement learning for the pick-up and delivery problem}{Nov. 2023 - Present}{}{}
% \item Developed an optimization-embedded deep reinforcement learning approach to solve the offline pick-up and delivery problems with service time windows.
% \item The approach outperforms the conventional model predictive control (MPC) algorithms and the state-of-the-art deterministic optimization solver (e.g., LKH3).
% \end{rSubsection}

\begin{rSubsection}{Vehicle routing problems in the semi-autonomous environment}{Jan. 2022 - June 2023}{}{}
\item Formulated a novel vehicle routing problem tailored to road networks with limited real-time support for autonomous vehicles.
\item Formulated a mixed-integer linear program (MILP) for the problem. Designed and implemented a two-phase tractable algorithm to effectively solve the problem of practical size.
\end{rSubsection}

% \begin{rSubsection}{Approximation Algorithms for Robust Shortest Path Problems}{Mar. 2022 - Present}{}{}
% \item Conducted a theoretical review of an NP-hard problem called the absolute robust shortest path problem (ARSP) and investigated its characteristics under various uncertainty sets of travel cost. 
% \item Devised an approximation algorithm to solve the ARSP under a convex hull uncertainty set of travel cost, with the goal of proving a constant approximation ratio of the algorithm.

% \end{rSubsection}
% \begin{rSubsection}{\href{https://publish.illinois.edu/cyber-risk/research-team-2/}{Cyber risk assessment for capital allocation}}{Jan. 2020 - Nov.2021}{}{}
% \item Proposed a data-verified stochastic loss model to capture the unique dynamics of cyber risk.
% \item Derived an optimal decision framework for risk managers to allocate ex-ante investment and ex-post loss reserve in cyber security.
% % \item Paper draft (R. Feng, W. Chong, L. Zhang, and H. Hu) submitted to \textit{\href{https://www.jri.pub/}{journal of risk and insurance}}.
% \end{rSubsection}



% \vspace{1 cm}
% \begin{rSubsection}{Next-generation multi-modal ride-sharing}{Oct. 2021 - April 2023}{}{}
% \item Designed an operational-level multi-modal ride-sharing system that seamlessly integrates mass transit and the first-and-last mile on-demand ride-sharing service.
% \item Devised an online optimization algorithm to generate nearly optimal strategies for request-vehicle matching, vehicle routing, and vehicle re-balancing, ensuring efficient operations of the system.
% % \item Conducted extensive simulations of the system with batched online requests sampled from a million taxi trips in NYC, providing valuable insights into the system's performance.
% % \item Paper draft (D. Edirimanna, H. Hu, S. Samaranayake) submitted to \textit{\href{https://aamas2022-conference.auckland.ac.nz}{AAMAS 2022}}
% \end{rSubsection}

% \begin{rSubsection}{\href{https://publish.illinois.edu/cyber-risk/research-team-2/}{Cyber risk assessment for capital allocation}}{Jan. 2020 - Nov.2021}{}{}
% \item Proposed a real-world-data-verified stochastic loss model to capture the unique dynamics of cyber risk.
% \item Derived an optimal decision framework for risk managers to allocate ex-ante investment and ex-post loss reserve in cyber security.
% % \item Paper draft (R. Feng, W. Chong, L. Zhang, and H. Hu) submitted to \textit{\href{https://www.jri.pub/}{journal of risk and insurance}}.
% \end{rSubsection}

% \begin{rSubsection}{Others}{}{}{}
% \item XXX
% \end{rSubsection}
\end{rSection}

\begin{rSection}{Selected Coursework}
% \vspace{-0.3cm}
\item \textbf{EECS:} Machine Learning, Deep Generative Models, Reinforcement Learning, Data Structure, Algorithm Design, Numerical Analysis, Simulation, Control Systems, Database Systems
\item \textbf{Optimization:} Mathematical Programming (Linear, Nonlinear, Mixed-Integer), Optimization Under Uncertainty, Combinatorial Optimization, Optimization of Logistics Systems
\end{rSection}

\begin{rSection}{Selected Publication}
\item \textbf{H. Hu}, et al. \textit{Column Generation for the Micro-Transit Zoning Problem}, Submitted to IJCAI 2026

\item S. Pradhan, B. Geary, \textbf{H. Hu}, et al., \textit{Hierarchical Reinforcement Learning for Designing Micro-Transit Zones}, Submitted to IJCAI 2026

\item \textbf{H. Hu}, R. Goswami, H. Jiang, S. Samarananyake, \textit{Optimal micro-transit zoning via clique
generation and integer programming}, IEEE ITSC 2025

\item H. Chen, Y. Wang, Y. Cai, \textbf{H. Hu}, J. Li, et al., \textit{HeuriGym: an agentic benchmark for LLM-crafted heuristics in combinatorial optimization}, NeurIPS 2025 MATH-AI workshop and ICLR 2026

\item Y. Kim, S. Li, \textbf{H. Hu}, W. Ouyang, S. Samaranayake, C. Wu, \textit{Learning to prune: fast feasible trip generation for high-capacity ridepooling}, TRISTAN 2025.

\item \textbf{H. Hu}, S. Samarananyake, \textit{Vehicle routing problems in the age of semi-autonomous driving}, In revision for the European Journal of Operational Research.

\item D. Edirimanna, \textbf{H. Hu}, S. Samarananyake, \textit{Integrating on-demand ride-sharing with mass transit at-scale}, In revision for Transportation Research Part C: Emerging Technologies.
\item M. Zalesak{$^*$}, \textbf{H. Hu}{$^*$}, S. Samarananyake, \textit{Ride-pool assignment algorithms: swapping heuristics and modern implementation}, Computers \& Operations Research. $^*$ \textit{Equal contributions}
\end{rSection}

\begin{rSection}{Skills}
Programming (Python, C++, MATLAB), Optimization Software (Gurobi), Machine Learning (PyTorch, PyG), Statistical Analysis (R), Database (SQL), Scientific Visualization
\end{rSection}

\begin{rSection}{Honors}
AWS Cloud Computing Grant (2024, Amazon), Graduate Student Fellowship (2021, Cornell), Summer Block Grant (2020, UIUC),  Scholarship for International Exchange Program (2018, CSC), Annual Academic Scholarship (2015-2018, SCUT)
\end{rSection}

\begin{rSection}{Professional Service}
I have served as reviewer for the following journals and conferences:
\begin{itemize}
    \renewcommand{\labelitemi}{-}
    \item The IEEE Intelligent Vehicles Symposium (IV 2026)
    \item The 40th Annual AAAI Conference on Artificial Intelligence (AAAI 2026)
    \item The Conference on Neural Information Processing Systems (NeurIPS 2025) 
    \item IEEE International Conference on Intelligent Transportation Systems (ITSC 2025)
    \item Learning for Dynamics \& Control Conference (L4DC 2024)
    \item Robotics: Science and Systems Conference (RSS 2024)
\end{itemize}
\end{rSection}

\begin{rSection}{Appointments}
\begin{rSubsection}{Teaching Assistant}{}{}{}
    \item Engineering Economics (CEE3230), Spring 2026
    \item Uncertainty Analysis in Engineering (CEE3040), Fall 2023, Fall 2024, Cornell
    \item \href{https://classes.cornell.edu/browse/roster/SP22/class/CEE/4665}{Modeling and Optimization for Smart Cities (CEE4665)}, Spring 2022, Cornell
    \item Probability and Statistics for Engineers (CEE202), Summer 2020 - Spring 2021, UIUC
\end{rSubsection}
\begin{rSubsection}{Research Assistant}{}{}{}
    \item US NSF/DOE Projects, Spring 2023, Spring 2024, Spring 2025, Cornell 
    \item\href{https://math.illinois.edu /illinois-risk-lab/home}{IRisk Lab},  Fall 2020 - Spring 2021, Department of Mathematics, UIUC
\end{rSubsection}
\end{rSection}

\end{document}



